%% kapa: exam=1 round=22 year=2011 period=2 form=A subject=appraisal_law subject_ko=부동산관계법규 seq=12 number=52
%% source: Q-Net 감정평가사 기출문제 (cst00309) artlSeq=5036790 file=extracted/1008/5036790/2011년 제22회 감정평가사 1차시험 기출문제/2교시.hwp
%% conversion: kordoc 4.24.51 -> tools/gen1.py (auto); review=proofread-gpt-astra; flags=none
\begin{question}{52}
국토의 계획 및 이용에 관한 법령상 개발행위허가에 관한 설명으로 옳지 않은 것은?
\begin{choices}
  \item 폐기물처리 사업자가 자연녹지지역안의 나대지에 폐지를 1개월 이상 쌓아 놓는 행위는 개발행위허가 대상이 아니다.
  % [원문 그대로] '제한 할'의 띄어쓰기 유지.
  \item 지구단위계획구역으로 지정된 지역으로서 도시관리계획상 특히 필요하다고 인정되는 지역에 대하여는 최고 5년 이내의 기간 동안 개발행위허가를 제한 할 수 있다.
  \item 개발행위허가권자는 공공시설의 귀속에 관한 사항이 포함된 개발행위허가를 하려면 미리 해당 공공시설이 속한 관리청의 의견을 들어야 한다.
  \item 특별시장ㆍ광역시장ㆍ시장 또는 군수는 개발행위허가를 하려면 해당 지역에서 시행되는 도시계획사업 시행자의 의견을 들어야 한다.
  \item 재해복구나 재난수습을 위한 응급조치는 개발행위허가를 받지 아니하고 할 수 있으나 1개월 이내에 특별시장ㆍ광역시장ㆍ시장 또는 군수에게 신고하여야 한다.
\end{choices}
\end{question}
\answer{1}
